% STATEMENT OF INTENT, warmer register, four paragraphs. Written 19 September.
%
% HIS BRIEF, and it is a deliberate departure from the two other versions in this
% folder: a letter checks boxes and nothing else. Show you know what you are
% getting into, and show you can do it. Read what the university and the programme
% claim as their strengths and link them explicitly to his interests. Read the
% prerequisites and link them to the CV and to experiences NOT on the CV. Four
% paragraphs: (1) introduce yourself and declare intent, (2) why you want to be
% here and how it benefits you, (3) why they benefit from having you, (4) thanks.
% Register warmer and more genuine than the other two, qualities shown through
% what he has actually done rather than asserted.
%
% HARD CAP OF 500 WORDS on the four paragraphs, set by him on 19 September, and
% the referees moved to a block at the foot so the body reads as prose rather
% than trailing off into administration. The referee block sits outside the 500.
%
% ONE THING HE WAS TOLD BEFORE IT WAS WRITTEN. The word "passion" and the
% declaration-of-enthusiasm register are the specific things the economics SOP
% guides he supplied on 18 September warn against, and that warning is recorded in
% CLAUDE.md. The warmth he asked for is here and it is carried by checkable facts:
% he built a full EU ETS research design for a thesis he then chose not to write,
% and he works through real analysis in his own time. Those say "I do this for
% fun" without the reader having to take his word for it. He can have the literal
% word if he wants it, and that is his call, not a future session's.
%
% LINDENBERGER IS OUT, his call on 19 September, and the reasoning is sound. The
% motivation had rested on the Kuemmel and Lindenberger argument that energy's
% output elasticity runs ahead of its cost share. That position is heterodox, it
% is a direct challenge to the cost-share theorem, and an auction theorist may
% well read it as fringe. Opening a letter on a contested claim spends the
% reader's goodwill on a fight nobody asked for. His own broader point survives
% without it: decarbonisation is a monumental effort, it reaches past capital into
% how societies work, and almost all of it runs through instruments governments
% design. That is mainstream, it is his, and it lands in the same place.
%
% WHAT CORE CLAIMS, verified 19 September off their own pages rather than inferred:
% founded 1966 by Jacques Drèze, modelled on the Cowles Foundation, "a globally
% renowned interdisciplinary research centre" bringing together "operations
% research, economics, and econometrics", more than 20 professors and 35 PhD
% students, and the training of young researchers at doctoral and postdoctoral
% stage named as a stated objective. The interdisciplinary claim is the one that
% does real work here, because his whole argument is that he is the econometrics
% half of a question the group approaches formally.
%
% PREREQUISITES LINKED, per his brief, including things not on the CV: Ockenfels
% on strategic thinking and competition at 1,3 (his field is market and auction
% design), mathematics 1,0, econometrics 1,3, numerical solutions to economic
% models, and independent real analysis. He confirmed the last two on 19 September.
% The thesis binned exposure, so the text says bins. Do not upgrade either claim.

% REVISION 19 September, third pass on this file.
%
% RELEVANCE AUDIT. Two passages were not earning their place and both are gone.
%   - "I am applying for the doctoral position on auction-based support schemes
%     for renewable energy" repeated the document's own header line for line.
%   - "founded in 1966 on that tradition, with more than twenty professors and
%     thirty-five doctoral students in one building" told Willems the age and
%     headcount of his own institute. The doctoral-student number survives, moved
%     to the one place it does work, which is the collegiality sentence.
%
% HEADINGS ADDED, his call on 19 September. CLAUDE.md had this recorded as a
% per-application judgement rather than an automatic yes, because on a one-page
% document headings cost lines. Here they fit: three small bold run-in headings,
% roughly four lines in total, paid for by the two cuts above.
%
% PERSONAL MATERIAL, from his own words on 19 September: he likes exchanging
% ideas with people from different cultures, backgrounds and disciplines, Erasmus
% at RSM Rotterdam showed him he works well in international environments, and he
% enjoys expressing what he cares about. Written into the Why CORE paragraph
% rather than bolted on, because at an institute that advertises an interdisciplinary
% remit and an international network it is a fit claim, not a personality note.
% "I enjoy expressing myself" is rendered as liking to explain what he is working
% on and hear what others are, which is seminar culture and is the research-group
% version of the same trait. The self-regarding version of that sentence would
% have cost him more than it gained.
%
% The Rotterdam semester was already on the CV and was doing nothing there. It is
% now evidence for something.

% REVISION 19 September, fourth pass, and this one is about joins rather than
% content. His objection: too many sentences stood alone. They were each true and
% none of them reached for the one before it, which is the same complaint he made
% about an early draft that was four disconnected blocks.
%
% THE THROUGH-LINE, stated so a future session does not flatten it again:
% decarbonisation is a rules problem before it is a technology problem; his
% training taught him to read what rules actually do, and the thesis is the proof
% because a policy looked inert until exposure was allowed for; that same blindness
% is what a support scheme poses; reading it properly needs theory he does not have
% and has started anyway; CORE is where the two halves sit in one building; he has
% worked alone and does not want to. Every paragraph hands the next one its
% subject.
%
% THE SENTENCE THAT DOES THE MOST WORK is in paragraph three: "What that taught me
% is that a policy can look inert until you allow for how differently places were
% exposed to it, which is the same problem a support scheme poses." Without it the
% thesis and the auction project are two blocks. With it they are one argument.
% Do not cut it for length.
%
% "The other half of what I want is people" is the hinge between the intellectual
% reason and the human one. It is short on purpose, between two long sentences.
%
% FRENCH B1 IS OUT, his call. Louvain-la-Neuve is francophone and B1 is enough to
% live there, but it was the last clause of a paragraph about econometrics and it
% read as a loose fact rather than an argument. It stays on the CV.

\documentclass[a4paper,11pt]{article}
\usepackage[utf8]{inputenc}
\usepackage[T1]{fontenc}
\usepackage[english]{babel}
\usepackage{lmodern}
\usepackage[left=23mm, right=23mm, top=17mm, bottom=15mm]{geometry}
\usepackage{parskip}
\usepackage{microtype}
\usepackage{xcolor}

\pagestyle{empty}
\setlength{\parskip}{0.45em}
\linespread{1.0}

\newcommand{\sh}[1]{\vspace{1.2mm}\noindent{\small\bfseries #1}\par\vspace{-0.15em}}

\begin{document}

\begin{center}
{\large\bfseries Statement of intent}\\[3pt]
{\small Baha Khammari \quad$\cdot$\quad PhD position, auction-based support schemes for renewable energy}\\
{\small LIDAM Energy Group, CORE, UCLouvain}
\end{center}

\vspace{3mm}

\sh{Why this doctorate}
Energy economics has followed me through both degrees at the University of Cologne, where I finish an M.Sc.\ this month with a thesis in applied microeconometrics. What kept me in it is that decarbonising our economies will be a monumental effort, reaching past capital into how our societies work, and that almost all of it runs through instruments governments design rather than technology alone. Which technologies get built, and at what cost to the public, is settled by those rules, and that is what I want to spend four years on.

\sh{Why CORE}
My training sits on the econometrics side of that question, and it has taken me to the edge of what econometrics alone answers. Reading an auction rule properly needs graduate industrial organization and game theory, which I do not have yet and have started on anyway, working through Jehle and Reny in my own time. It is why I am writing to CORE rather than a department that does one of these things, because a centre built around operations research, economics and econometrics is where both halves of the question already sit. The other half of what I want is people. Everything I have written so far I wrote alone, and an Erasmus semester in Rotterdam showed me I work better in a room where nobody shares my background, explaining what I am working on and hearing what others are. That is most of what a research group does, and I want it among your doctoral students and visiting researchers.

\sh{What I would bring}
My thesis, graded 1.0 and worth 30 of 120 ECTS, asked whether a national scholarship programme in Brazil raised formal-sector wages, on a panel I built from administrative records, five million linked employer-employee observations across 558 microregions, 2005 to 2019, in R through BigQuery. Callaway and Sant'Anna (2021) group-time effects with doubly-robust inference across exposure bins put the gain at 3.04 per cent in low-exposure regions, distinguishable from zero at the 10 per cent level, on a non-monotonic dose-response whose sign two-way fixed effects reverses. What that taught me is that a policy can look inert until you allow for how differently places were exposed to it, which is the same problem a support scheme poses. Since the 2016 state aid framework member states have arrived at technology-neutral, technology-specific and mixed designs at different dates, so the variation is staggered adoption of a design rather than of a policy, and the AURES II database records every renewable auction in the EU and UK from 2011 to 2021. I would use it to estimate how joint bidding shifts the awarded portfolio and the rent per megawatt hour against member states not yet treated. Behind the econometrics sit advanced microeconomics on strategic thinking and competition with Axel Ockenfels at 1.3, 1.0 in mathematics, and a model I calibrated in MATLAB.

Thank you for considering my application. I could start in October 2026 and would be glad to talk the project through whenever it suits you.

\vspace{3.5mm}
{\color{black!35}\hrule height 0.5pt}
\vspace{2.5mm}

\begin{flushleft}
{\small\textbf{Referees}}\\[3pt]
{\small Prof.\ Dr.\ Pia Pinger, University of Cologne, chair of my master's thesis.\\
\texttt{pia.pinger@uni-koeln.de}}\\[4pt]
{\small Jun.-Prof.\ Dr.\ Oliver Ruhnau, University of Cologne and EWI, examiner in \textit{Energy Markets and Regulation}.\\
\texttt{ruhnau@wiso.uni-koeln.de}}
\end{flushleft}

\vspace{2mm}
\begin{flushleft}
{\small Baha Khammari \quad$\cdot$\quad \texttt{b.e.khammari@gmail.com} \quad$\cdot$\quad +49\,1520\,9016956}
\end{flushleft}

\end{document}
